\section*{Meaning of the notation}

The papers reviewed below reuse some letters for different internal
quantities. To make the comparison unambiguous, this chapter adopts the common
notation in Table~\ref{tab:notation}. Every displayed equation is restated with
this notation, even when the cited paper uses different letters. For temporal
comparisons, $\tau$ is the common forecast origin (the latest observed sample)
and $d$ is a backward offset from that origin. Setting $\tau=0$ therefore puts
the latest observed value at relative time $0$, the preceding values at
$-1,-2,\ldots$, and the first forecast at $+1$. The paper-specific rows in the
table translate each original convention to this common axis. The other index
roles remain fixed throughout: $k$ denotes a target series, $j$ a source
series, $h$ a forecast horizon, $i$ an explained SHAP component, $n$ a test
example, and $r$ a training run. Superscripts distinguish quantities that the
original papers denote by the same Greek letter; for example, attention
coefficients and DCIts transition coefficients are written as
$\alpha^{\mathrm{att}}$ and $\alpha^{\mathrm{DCI}}$, respectively. These
superscripts clarify the cross-paper comparison without changing the
underlying definitions.

\begingroup
\small
\setlength{\LTpre}{0.4\baselineskip}
\setlength{\LTpost}{0.7\baselineskip}
\renewcommand{\arraystretch}{1.08}
\begin{longtable}{@{}
>{\raggedright\arraybackslash}p{0.24\textwidth}
>{\raggedright\arraybackslash}p{0.69\textwidth}@{}}
\caption{Common notation and time-origin conventions used across the reviewed
benchmarks.}
\label{tab:notation}\\
\toprule
Symbol & Meaning \\
\midrule
\endfirsthead
\multicolumn{2}{@{}l}{\small\emph{Table~\thetable\ continued}}\\
\toprule
Symbol & Meaning \\
\midrule
\endhead
\midrule
\multicolumn{2}{r@{}}{\small\emph{Continued on the next page}}\\
\endfoot
\bottomrule
\endlastfoot

$\tau$, $h$, $H$ &
$\tau$ is the common forecast origin and latest observed time;
$h\in\{1,\ldots,H\}$ is a forecast step; and $H$ is the complete forecast
horizon. The forecasted time is $\tau+h$. On the relative axis $\tau=0$, the
first unseen target is $h=1$, at time $+1$. \\

$d$, $L$ &
$d\in\{0,\ldots,L-1\}$ is the backward offset from the common origin and $L$
is the input-window length. Thus $x_{j,\tau-d}$ denotes a historical input,
with $d=0$ denoting the latest observed sample. \\

$t$, $\ell$ &
$t$ and $\ell$ reproduce the indexing used in a paper's own equation. In the
target-time form $x_{j,t-\ell}\rightarrow x_{k,t}$, $\ell=1$ is the immediately
preceding sample. For one-step prediction, $t=\tau+1$ and the common offset is
$d=\ell-1$. Paper-specific uses of $t$ are listed below. \\

$j$, $k$, $N$ &
$j$ indexes a source or input series, $k$ indexes a target series, and $N$ is
the number of series. \\

$i$, $P$ &
$i$ indexes one of the $P$ feature groups or components treated as players in
the SHAP calculation. \\

$n$, $m$; $r$, $R$ &
$n$ indexes one of $m$ explained test examples; $r$ indexes one of $R$
repeated training runs. \\

$x$, $x^{(n)}$, $X$ &
$x$ is a fixed forecasting input, $x^{(n)}$ is test example $n$, and $X$ is a
random replacement input drawn by a background or dependency-aware sampler. \\

$x_{j,\tau-d}$; $x_{j,t-\ell}$; $x_{k,t}$; $y_t$ &
$x_{j,\tau-d}$ uses the common origin convention. $x_{j,t-\ell}$ uses the
target-time convention in the Bari\'c and MDA generator equations;
$x_{k,t}$ is target series $k$ at time $t$; and $y_t$ denotes a scalar target
when a paper separates the target from its input series. \\

$u_n$, $d_{\mathrm{rec}}$ &
$u_n$ is term $n$ of a recurrent sequence and $d_{\mathrm{rec}}$ is the
effective recurrence degree: the largest backward dependency allowed by the
sampled expression. The subscript ``rec'' distinguishes recurrence degree from
the common forecast-origin offset $d$. \\

$\mathcal{T}$, $o$, $F_{\mathcal{T}}$ &
$\mathcal{T}$ is a sampled unary--binary expression tree, $o$ is its number of
operator nodes, and $F_{\mathcal{T}}$ is the recurrence function obtained by
executing that tree. \\

$\mathbf z$, $\mathbf h_v$, $\boldsymbol\pi_v$, $s_v$, $a(s_v)$ &
$\mathbf z$ is a continuous latent code for an expression tree;
$\mathbf h_v$ is the decoder state at tree node $v$;
$\boldsymbol\pi_v$ is the predicted distribution over the symbol library;
$s_v$ is the selected symbol; and $a(s_v)\in\{0,1,2\}$ is its arity. These
symbols describe the hierarchical variational autoencoder of
\citet{meznar2023efficient}. \\

$e$, $\mathcal D$, $\mathcal H$; $\mathbf z_{\mathrm{num}}$,
$\mathbf z_{\mathrm{hyp}}$, $\mathbf z_{\mathrm{fuse}}$ &
$e$ is a symbolic expression; $\mathcal D=\{(\mathbf x_i,y_i)\}_{i=1}^{n}$ is
its numerical support; and $\mathcal H$ is a tokenized set of structural
hypotheses. In NSRwH, the numerical and hypothesis encoders produce
$\mathbf z_{\mathrm{num}}$ and $\mathbf z_{\mathrm{hyp}}$, whose elementwise
sum $\mathbf z_{\mathrm{fuse}}$ conditions expression decoding. \\

$D_{\mathrm{ode}}$, $\mathbf x(t)$, $\mathbf f$, $\widehat{\mathbf f}$,
$\Phi_{\mathbf f}$ &
$D_{\mathrm{ode}}$ is the number of state variables in an ODE system;
$\mathbf x(t)\in\mathbb R^{D_{\mathrm{ode}}}$ is its state;
$\dot{\mathbf x}=\mathbf f(\mathbf x)$ is the sampled vector field; and
$\widehat{\mathbf f}$ is the equation recovered from observations. The flow
$\Phi_{\mathbf f}(\Delta;\mathbf x_0)$ is the state reached after elapsed time
$\Delta$ by integrating $\mathbf f$ from initial condition $\mathbf x_0$. \\

$M_{\mathrm{S^2}}$, $N_{\mathrm{S^2}}$, $T_{\mathrm{S^2}}$;
$\mathbf X$, $\mathbf Y$, $\mathbf f_{\mathrm{S^2}}$ &
$M_{\mathrm{S^2}}$ and $N_{\mathrm{S^2}}$ are the numbers of input and output
channels in S$^2$Generator; $T_{\mathrm{S^2}}$ is the generated sequence
length; $\mathbf X\in\mathbb R^{M_{\mathrm{S^2}}\times T_{\mathrm{S^2}}}$
is the sampled temporal input; and
$\mathbf Y=\mathbf f_{\mathrm{S^2}}(\mathbf X)\in
\mathbb R^{N_{\mathrm{S^2}}\times T_{\mathrm{S^2}}}$ is the pointwise
symbolic response. The S$^2$ subscripts prevent collision with the common
series count $N$ and input-window length $L$. \\

$\mathcal K$, $J_K$, $R_K$, $\kappa_r$, $\kappa_\star$ &
$\mathcal K$ is a bank of Gaussian-process covariance kernels; $J_K$ is the
maximum number selected for one synthetic series; $R_K$ is the sampled number;
$\kappa_r$ is selected kernel $r$; and $\kappa_\star$ is their composite
kernel. These symbols are used for KernelSynth and its descendants. \\

$T_{\mathrm{syn}}$, $N_{\mathrm{lat}}$, $z_a(t)$,
$\lambda_{j,a}$ &
$T_{\mathrm{syn}}$ is the generated-series length; $N_{\mathrm{lat}}$ is the
number of latent functions in LMC-Synth; $z_a(t)$ is latent function $a$; and
$\lambda_{j,a}$ is its nonnegative mixing weight in observed channel $j$,
with $\sum_a\lambda_{j,a}=1$. \\

$P_{\mathrm{patch}}$, $D_{\mathrm{model}}$ &
$P_{\mathrm{patch}}$ is the number of consecutive observations represented by
one model token and $D_{\mathrm{model}}$ is the token-embedding dimension.
The subscript distinguishes patch length from the $P$ SHAP feature groups
defined above. \\

$\mathbf x_C$, $\mathbf K_{CC}$, $\mathbf k_{hC}$,
$\mu_{h\mid C}$ &
$\mathbf x_C$ is an observed Gaussian-process context; $\mathbf K_{CC}$ is
its covariance matrix under $\kappa_\star$; $\mathbf k_{hC}$ contains the
covariances between future step $h$ and the context; and $\mu_{h\mid C}$ is
the conditional predictive mean at that step. \\

\addlinespace
\multicolumn{2}{@{}l}{\textit{Time-origin convention in each reviewed paper}}\\
\addlinespace

SHAPformer (2026) &
The task contains 168 observed hours followed by 168 forecast hours. In the
authors' generator arrays, past positions are $0,\ldots,167$ and future
positions are $168,\ldots,335$. Re-centering at raw position 167 gives
$\tau=0$: raw position 167 is the latest observation, raw position 168 is
$h=1$, and output-array position 0 also denotes $h=1$. \\

Bari\'c et al. (2021) &
The generator is written at the target time:
$x_{k,t}\leftarrow x_{j,t-\ell}$. Hence $t-1$ is the latest available sample
for predicting $t$. Under the common origin, $\tau=t-1$, so that sample is
relative time 0 and the target $t$ is relative time $+1$; paper lag $\ell$
maps to $d=\ell-1$. Zero-based heatmap columns are array positions rather than
signed relative times. \\

D'Ascoli et al. (2022) &
The recurrence target is $u_n$ and its admissible historical leaves are
$u_{n-i}$ for $i\in\{1,\ldots,d_{\mathrm{rec}}\}$. For a one-step forecast of
$u_n$, the common origin is $\tau=n-1$: $u_{n-1}$ has relative position 0,
$u_{n-2}$ has position $-1$, and $u_n$ is at $h=1$. \\

Me\v{z}nar et al. (2023) &
The paper's coordinate is the static input vector in expressions such as
$y=f(x_1,\ldots,x_p)$. When adapted to this thesis, a leaf
$\operatorname{LAG}(j,\ell)$ follows the
common convention: $\ell=1$ is the latest value available for a one-step
forecast and maps to $d=0$. \\

Bendinelli et al. (2023) &
NSRwH also learns static expressions $y=e(\mathbf x)$ from unordered numerical
input--output pairs. A time-series adaptation would construct each row of
$\mathbf x$ from lagged observations; under the common convention, a feature
$\operatorname{LAG}(j,1)$ contains the latest observation at $d=0$. \\

ODEFormer (2024) &
The observations have continuous timestamps $t_i$. After selecting an origin
$t_0$, the initial state $\mathbf x(t_0)$ lies at elapsed time $\Delta=0$ and
future states lie at $\Delta=t-t_0>0$. The sampled tree defines the
instantaneous derivative $\dot{\mathbf x}(t)$ at the current state; numerical
integration converts that local law into an entire future trajectory. \\

MDA (2024) &
The synthetic equations also use target-time indexing:
$y_t\leftarrow x_{j,t-\ell}$. Therefore the latest input is $t-1$ and becomes
$\tau=0$ after setting $\tau=t-1$. The paper labels the first forecast as
$h=1$; the implementation stores it at output-array index 0. \\

CrossScaleNet (2025) &
The input is written chronologically as $[x_1,\ldots,x_T]$, and the saliency
figures use zero-based window positions $0,\ldots,L-1$: position 0 is the
oldest displayed input and position $L-1$ the latest, which becomes $\tau=0$.
Table~1 calls small important-lag numbers ``recent''. The paper and repository
leave the exact synthetic generator equation unpublished, so the conversion
between each lag label and each plotted column remains unspecified. \\

S$^2$Generator and SymTime (2025) &
The symbolic stage uses aligned timestamps:
$y_{i,t}=f_i(x_{1,t},\ldots,x_{M_{\mathrm{S^2}},t})$. If an aligned sample is
chosen as the origin, $\mathbf x_{\tau}$ and $\mathbf y_{\tau}$ both have
relative time 0. Temporal lags belong to the separately sampled ARMA process
that can generate an input channel $x_j$; the leaves of $f_i$ refer to current
inputs $x_{j,t}$. \\

DCIts (2026) &
The paper anchors time at the end of the observed window:
$Q_t=(X_t,X_{t-1},\ldots,X_{t-L+1})$ and predicts $X_{t+1}$. Thus the latest
sample is already $t=\tau$ (relative time 0), and the target is $t+1$
($h=1$). With the paper's one-based lag coordinate,
$(Q_t)_{j,\ell}=x_{j,t-\ell+1}$; hence $\ell=1$ is the latest sample and maps
to $d=0$. \\

$g=(g_1,\ldots,g_H)$; $f$; $\widehat g$ &
$g$ is the known synthetic generator and $g_h(x)$ its output at horizon $h$;
$f$ is a forecaster learned from data generated by $g$; and $\widehat g$ is a
symbolic mechanism identified independently from the generated observations. \\

$S$, $\bar S$, $q(\cdot\mid x_S)$ &
$S\subseteq\{1,\ldots,P\}$ is a coalition of retained components, $\bar S$ is
its complement, and $q$ is the dependency-aware distribution used to resample
the absent components conditional on the retained values. \\

$v_{g,h}(S;x)$; $\phi_{i,h}(x)$; $G_i$ &
$v_{g,h}$ is the expected generator output for coalition $S$ at horizon $h$;
$\phi_{i,h}$ is the signed local SHAP contribution of component $i$; $G_i$ is
its global absolute-importance percentage. \\

$M^{\mathrm{GT}}_{k,j,\ell}$; $B^{\mathrm{GT}}_{k,j}$ &
Binary generator-derived references: $M^{\mathrm{GT}}$ marks active
target--source--lag entries and $B^{\mathrm{GT}}$ marks target--source edges
after collapsing the lag axis. \\

$\mathbb{I}[\cdot]$ &
Indicator function: it equals 1 when the statement in brackets is true and 0
otherwise. \\

$b_k$, $c_{k,j,\ell}$, $\varepsilon_{k,t}$ or $\varepsilon_t$ &
Generator bias, signed source--lag coefficient, and innovation or noise; the
target index is omitted for a scalar target. \\

$\odot$ &
Elementwise (Hadamard) multiplication of arrays with compatible shapes. \\

$\widehat{\mathcal E}_k$, $\widehat\ell_{k,j}$ &
TCDF's retrieved directed source--target edges for target $k$ and the retrieved
delay from source $j$. \\

$\alpha^{\mathrm{att}}$, $\beta^{\mathrm{att}}$ &
Temporal and source-variable attention quantities, respectively, in the
seq2graph/IMV-LSTM discussion; both are nonnegative normalized weights. \\

$\alpha^{\mathrm{va}}$, $\alpha^{\mathrm{ta}}$,
$\alpha^{\mathrm{c}}$ &
MDA variable-attention, temporal-attention, and combined feature--time scores;
$\alpha^{\mathrm{c}}=\alpha^{\mathrm{va}}\odot\alpha^{\mathrm{ta}}$. \\

$Q_t$, $F_t$, $C_t$, $\alpha^{\mathrm{DCI}}_t$;
$(\alpha^{\mathrm{DCI}}_t)_{k,j,\ell}$ &
DCIts input window, Focuser mask, signed Modeler tensor, and final signed
transition tensor, with $\alpha^{\mathrm{DCI}}_t=C_t\odot F_t$; the indexed
entry links source $j$ at lag $\ell$ to target $k$. In the paper's ordering,
$(Q_t)_{j,\ell}=x_{j,t-\ell+1}$, so $\ell=1$ contains the latest observation
$x_{j,t}$. \\

$\bar\alpha^{\mathrm{DCI},(r)}_{k,j,\ell}$ &
DCIts coefficient averaged over held-out windows within training run $r$; the
bar and run superscript make the two aggregation levels explicit. \\

$\mu_{k,j,\ell}$, $\sigma_{k,j,\ell}$,
$I^{\mathrm{stable}}_{k,j,\ell}$ &
Across-run mean, across-run standard deviation, and DCIts stability indicator
for a reported transition coefficient. \\

$\widetilde\beta^{\mathrm{DCI}}_{t,k,j}$,
$\beta^{\mathrm{DCI}}_{t,k,j}$ &
DCIts local absolute coefficient strength aggregated over lags and its
source-wise normalized form; aggregation removes lag and sign. The index $t$
is omitted when the same calculation is applied to a coefficient tensor that
has already been averaged over held-out windows. \\

$\mathcal M(g)$, $\mathcal M(\widehat g)$ &
Comparable representations of the true and identified mechanisms, such as
their source--target--lag support, signed coefficients, regimes, or symbolic
terms. \\

$E_g(x)$, $E_f(x)$, $E_{\widehat g}(x)$, $A(f,x)$ &
Explanation of the known generator, model-level reference explanation of the
trained forecaster, explanation of the independently identified symbolic
model, and explanation estimated for $f$ by an XAI method, respectively. \\

$d_y$, $d_{\mathcal M}$, $d_{\mathcal E}$ &
Task-appropriate distances for outputs, mechanism representations, and
explanations. Their concrete forms depend on the reference family: for
example, MSE for forecasts, precision/recall for structural support, coefficient
error for signed mechanisms, or attribution error and rank agreement for local
explanations. \\

\end{longtable}
\endgroup

